\subsubsection{Hardware Revisions}
The pin misassignments identified in this revision will be corrected in a dedicated PCB revision, restoring interrupt-driven edge detection for the Gate 03 input and  enabling PWM-driven LED control for the slide potentiometer indicators. 

Additionally, the fourth WS2812 RGB LED, currently unmapped in firmware, will be assigned and integrated into the LED control scheme.
An aluminium front panel with UV-printed graphics is planned to replace the current  prototype panel, bringing the module to a finished Eurorack form factor.

\subsubsection{Memory and Buffer Capacity}
To address the current memory constraints and extend the recording buffer beyond  \qty{2.5}{\second}, an external RAM expansion via QSPI is planned. 
This would also provide the headroom necessary to accommodate additional DSP algorithms without further pressure on the internal SRAM.

A suitable candidate for this extension is the APS6404L-3SQR-SN \cite{APS6404L_3SQR}, a 64 Mbit (8 MB) PSRAM in an SOP-8 package.

At 8 MB, the chip would increase available RAM by approximately eight times thecurrent internal SRAM, enabling significantly longer recording buffers and providing headroom for future DSP algorithm development.

\subsubsection{DSP Algorithms}
The current DSP chain will be extended with additional audio processing algorithms. 
Planned additions include a synced delay effect, filtering, and grain synthesis or sampling functionality, expanding the module's capabilities.

%\subsubsection{Noise Characterisation}
%A formal noise floor measurement using the swept-sine method described in  Section~\ref{sec:frequency-response} is planned for a subsequent revision, to formally verify the effectiveness of the analogue/digital separation strategy applied during PCB layout.